% \iffalse meta-comment
% 
% File: tipauni.dtx
% ---------------------------------------------------------------------------
% Package:      tipauni
% Author:       निरंजन
% Version:      v0.6 (21 February, 2022)
% Description:  For producing Unicode characters with TIPA commands.
% Repository:   https://git.gnu.org.ua/tipauni.git
% Bug tracker:  https://puszcza.gnu.org.ua/bugs/?group=tipauni
% License:      GPL v3.0+, GFDL v1.3+
% ---------------------------------------------------------------------------
% LaTeX Package tipauni v0.6
% Copyright © 2021, 2022 निरंजन
% 
% This program is free software: you can redistribute it and/or modify
% it under the terms of the GNU General Public License as published by
% the Free Software Foundation, either version 3 of the License, or
% (at your option) any later version.
% 
% This program is distributed in the hope that it will be useful,
% but WITHOUT ANY WARRANTY; without even the implied warranty of
% MERCHANTABILITY or FITNESS FOR A PARTICULAR PURPOSE. See the
% GNU General Public License for more details.
% 
% You should have received a copy of the GNU General Public License
% along with this program. If not, see <https://www.gnu.org/licenses/>.
% 
% The current maintainer of this work is निरंजन.
% 
% This work consists of the files tipauni.dtx
%                                 tipauni.ins
% and the derived file            tipauni.sty.
% \fi
% \iffalse
%<*internal>
\iffalse
%</internal>
%<*readme>
--------------------------------------------------------------------------
Package:      tipauni
Author:       निरंजन
Version:      v0.6 (21 February, 2022)
Description:  For producing Unicode characters with TIPA commands.
Repository:   https://git.gnu.org.ua/tipauni.git
Bug tracker:  https://puszcza.gnu.org.ua/bugs/?group=tipauni
License:      GPL v3.0+, GFDL v1.3+
--------------------------------------------------------------------------
%</readme>
%<*internal>
\fi
%</internal>
%<*driver|package>
\def\tipauniname{tipauni}
\def\tipauniversion{0.6}
\def\tipaunidate{2022-02-21}
\def\tipaunidescription{For producing Unicode characters with TIPA commands}
%</driver|package>
%<*driver>
\documentclass{l3doc}
\usepackage{multicol}
\usepackage{graphicx}
\usepackage{array}
\usepackage{longtable}
\usepackage[referable]{threeparttablex}
\usepackage{booktabs}
\usepackage{float}
\usepackage{xcolor}
\usepackage{hologo}
\DeclareRobustCommand\XeLaTeX{\hologo{XeLaTeX}}
\usepackage{fontspec}
\newfontfamily{\ipa}[Color=blue]{CharisSIL}
\newfontfamily{\devtext}[%
Script=Devanagari,%
Renderer=Harfbuzz,%
Scale=0.8%
]{Shobhika}
\usepackage{hyperref}
\hypersetup{%
  colorlinks,%
  linkcolor=red!60!black,%
  citecolor=green!60!black,%
  urlcolor=blue!60!black%
}%
\usepackage{fontawesome5}
\renewcommand{\arraystretch}{1.2}

\begin{document}
\DocInput{\tipauniname.dtx}
\end{document}
%</driver>
% \fi
% \title{The \pkg{\tipauniname} package}
% \author{^^A
%   {\devtext निरंजन}^^A
%   \thanks{^^A
%     \faIcon{envelope}\quad
%     \href{mailto:hi.niranjan@pm.me}^^A
%     {\ttfamily hi.niranjan@pm.me}^^A
%   }^^A
% }
% \date^^A
% {^^A
%   Version \tipauniversion\ --- \tipaunidate\\[1ex]^^A
%   {^^A
%     \small\faIcon{link}\quad
%     \url{https://ctan.org/pkg/tipauni}^^A
%   }\\^^A
%   {^^A
%     \small\faIcon{bug}\quad
%     \url{https://puszcza.gnu.org.ua/bugs/?group=tipauni}^^A
%   }^^A
% }^^A
% \maketitle
%   
% \tableofcontents
%   
% \vfill
% {%
%   \setlength{\parindent}{0pt}
%   Copyright © 2021, 2022 {\devtext निरंजन}
%   
%   Permission is granted to copy, distribute and/or modify this document under
%   the terms of the GNU Free Documentation License, Version 1.3 or any later
%   version published by the Free Software Foundation; with no Invariant
%   Sections, no Front-Cover Texts, and no Back-Cover Texts. A copy of the
%   license is included in the section entitled ``GNU Free Documentation
%   License''.
% }%
%
% \clearpage\pagebreak
% 
% \begin{documentation}
%   \begin{abstract}
%     This package is the second step for \LaTeX-internally bridging the gap
%     between the Unicode and the TIPA package. The first step was the package
%     \href{https://www.ctan.org/pkg/unitipa}{\pkg{unitipa}}. Users can now use
%     their beloved TIPA shortcuts with the benefits of Unicode i.e.\
%     searchability, copy-pasting, changing the font and many more. As this
%     package needs the package \pkg{fontspec} for loading an IPA font, it needs
%     to be compiled with Lua/\XeLaTeX. This package can also be viewed as an
%     ASCII-based input method for producing IPA characters in Unicode.
%   \end{abstract}
%   
%   \section{Acknowledgements}
%   
%   This package is completed because of the rigorous help by Jairo A.\ del Rio,
%   Jonathan P.\ Spratte, Phelype Oleinink \textit{\&} Ulrike Fischer. Their
%   suggestions and contributions are extremely valuable for the development of
%   this package.
%   
%   Jonathan P.\ Spratte has helped the package by building a parser based on
%   \pkg{expl3}. Phelype Oleinik has helped the package by building a test-suit
%   for regression-tests. The tests can be seen in the git-repository of the
%   project.
%
%   Lemures Lemniscati pointed a few bugs in the documentation \& the code.
%   Their details are as follows.
%   are as follows:
%   
%   \begin{description}
%   \item[2022-01-03]
%     A typo '\cs{textsh}' in tipauni.dtx.
%     \href{https://puszcza.gnu.org.ua/bugs/index.php?538}{(link)}
%   \item[2022-01-09]
%     Avoid unconditional loading of CharisSIL.
%     \href{https://puszcza.gnu.org.ua/bugs/index.php?539}{(link)}
%   \item[2022-01-09]
%     Fix typos in Input B of Raised and Lowered.
%     \href{https://puszcza.gnu.org.ua/bugs/index.php?540}{(link)}
%   \item[2022-01-09]
%     \cs{textcorner} should be defined by \cs{DeclareTextSymbol}.
%     \href{https://puszcza.gnu.org.ua/bugs/index.php?541}{(link)}
%   \item[2022-01-09]
%     Fix a typo in Input B of Near-close near-front rounded.
%     \href{https://puszcza.gnu.org.ua/bugs/index.php?542}{(link)}
%   \end{description}
%   
%   All of these contributions are highly valuable for the development of this
%   package.
%   
%   \section{Package options}
%   
%   \begin{function}{documentfont}
%     \begin{syntax}
%       documentfont=\meta{font-name}
%     \end{syntax}
%     
%     The default font for the entire document set by this package is the
%     \href{https://ctan.org/pkg/charissil}{Charis SIL}. Users can easily change
%     it with the package-option \verb|documentfont|. This option takes one
%     value which is the font-name. e.g.\ For using the
%     \href{https://ctan.org/pkg/doulossil}{Doulos  SIL} font as the main font,
%     use the package like it is used in the following.
%     
% \begin{verbatim}
% \usepackage[documentfont=DoulosSIL]{tipauni}
% \end{verbatim}
%     
%     Please note that you must provide a value to this key. If you pass it an
%     empty value; i.e.,
%     \enquote{\verb|documentfont=\{\}|}/\enquote{\verb|documentfont=|} or no
%     value, i.e., \enquote{\verb|documentfont|}, you will be prompted with
%     errors from package \pkg{tipauni}. If an empty/no-value
%     \verb|documentfont| is used \& no other non-empty \verb|documentfont| has
%     been used so far; then the document is typeset with the default
%     font-family of Lua/\XeLaTeX, i.e., the Latin Modern family. If another
%     \verb|documentfont| has a non-empty value, the value of that option will
%     be used as the font of the document.
%   \end{function}
%   
%   \begin{function}{fontspecoptions}
%     \begin{syntax}
%       fontspecoptions=\meta{font-options}
%     \end{syntax}
%     
%     This option can be used to set options to the font set with
%     \verb|documentfont| package option. If the \verb|recommendedfont| option
%     is used \& has the highest priority; the options set with
%     \verb|fontspecoptions| will be used with the Charis SIL font. If the
%     \verb|preservefont| option has the highest priority, this parameter will
%     be ineffective \& throw a warning. In the argument of this option write as
%     if you are writing in the optional parameter of the \cs{setmainfont}
%     command. e.g.\
%     
% \begin{verbatim}
% \usepackage[fontspecoptions={Scale=1.5}]{tipauni}
% \end{verbatim}
%     
%     \noindent If this options is used with no value it will be ignored with a
%     warning.
%     
%     \noindent e.g.\ Try \verb|\usepackage[fontspecoptions]{tipauni}|.
%     
%     \medskip
%
%     \noindent If it is used with an empty value it will reset the fontspec-options so
%     far you have set. e.g.\ Try:
%     
% \begin{verbatim}
% \usepackage[%
%   fontspecoptions={Scale=1.5},%
%   fontspecfontspecoptions={}%
% ]{tipauni}
% \end{verbatim}
%   \end{function}
%   
%   Please note that the keys so far documented are the only keys which take
%   values. Other than these the following keys \emph{don't} take any value. You
%   will be prompted with an error from package \pkg{expkv} if you use values
%   for the following keys.
%   
%   \begin{function}{preservefont}
%     This option will omit executing the \cs{setmainfont} command with Charis
%     SIL. The typeset document will use the default font of Lua/\XeLaTeX, i.e.,
%     the Latin Modern.
%   \end{function}
%   
%   \begin{function}{recommendedfont}
%     This option is used to override other font definitions made with option
%     \verb|documentfont|. Charis SIL will be used as the main font of the
%     document when this option is used.
%   \end{function}
%   
%   Please note that every option so far mentioned can potentially override each
%   other. e.g.\ If a \verb|documentfont| is added after a \verb|preservefont|,
%   the former will override the latter (\& vice versa). It is true for all
%   these three package options.
%   
%   \begin{function}{resetfontspecoptions}
%     This option is used to reset all the \pkg{fontspec} options set with
%     package option \verb|fontspecoptions|. As noted for the font-options,
%     these two options which deal with the options of the loaded fonts also can
%     override each other. The last one loaded will be considered of the highest
%     priority.
%   \end{function}
%   
%   \begin{function}{incompatible}
%     By default \pkg{\tipauniname} will throw a warning if symbols not
%     contained in the \pkg{tipa} package were used. Using this option you can
%     switch this check off if you don't have to stay compatible with
%     \pkg{tipa} to silence this warning.
%   \end{function}
%   
%   \section{List of available commands}
%   TIPA has two alternative methods for inputting IPA characters. In the
%   following tables, \textbf{Input A} column lists all the long TIPA commands
%   (e.g.~\verb|\textrtailt|) \textit{\&} \textbf{Input B} column lists the
%   short ones (e.g.~\verb|\textipa{\:t}|)\footnote{All of these commands must
%   be enclosed in the \texttt{\textbackslash textipa}
%   command. (e.g.~\texttt{\textbackslash textipa\{\textbackslash :tP\}} for
%   {\ipa [ʈʔ]})}. This documentation lists down all the commands necessary for
%   typing symbols in the IPA chart.
%   
%   \subsection{Pulmonic consonants}
%   
%   \begin{ThreePartTable}
%     \begin{TableNotes}
%     \item[\textcolor{red}{§}] These commands are also available without
%       package \pkg{tipa}.
%     \item[1] \label{1} A symbol for this sound was introduced in 2005. The
%       last update of the package \pkg{tipa} was in 2004, so this character
%       is missing in it. A command for it is provided by this package. Such
%       \pkg{tipauni}-exclusive commands are written with red color in this
%       documentation.
%     \end{TableNotes}
%     \small
%     \begin{longtable}{%
%       l@{\hspace{0.3cm}-\hspace{0.3cm}}%
%       l@{\hspace{0.3cm}-\hspace{0.3cm}}%
%       l@{\hspace{0.3cm}-\hspace{0.3cm}}%
%       >{\ipa[}l<{]}%
%     }%
%       \toprule
%       \endfirsthead
%       \toprule
%       \endhead
%       \bottomrule
%       \endfoot
%       Description & Input A & Input B &
%       \multicolumn{1}{l}{\kern-1.25ex Symbol}\\
%       \midrule
%       Retroflex voiceless stop & \verb|\textrtailt| & \verb|\:t| & ʈ \\
%       Retroflex voiced stop & \verb|\textrtaild| & \verb|\:d| & ɖ \\
%       Palatal voiced stop & \verb|\textbardotlessj| & \verb|\*j| & ɟ \\
%       Uvular voiced stop & \verb|\textscg| & \verb|\;G| & ɢ \\
%       Glottal voiceless stop & \verb|\textglotstop| & \verb|P| & ʔ \\
%       Labiodental nasal & \verb|\textltailm| & \verb|M| & ɱ \\
%       Retroflex nasal & \verb|\textrtailn| & \verb|\:n| & ɳ \\
%       Palatal nasal & \verb|\textltailn| & \verb|\*n| & ɲ \\
%       Velar nasal & \verb|\ng|\tnote{\textcolor{red}{§}} & \verb|N| & ŋ \\
%       Uvular nasal & \verb|\textscn| & \verb|\;N| & ɴ \\
%       Bilabial trill & \verb|\textscb| & \verb|\;B| & ʙ \\
%       Uvular trill & \verb|\textscr| & \verb|\;R| & ʀ \\
%       \multicolumn{1}{>{\color{red}}l@{\hspace{0.3cm}-\hspace{0.3cm}}}{Labiodental flap\tnotex{1}} &
%       \color{red}\verb|\labdentflap| & \color{red}\verb|\*v| & ⱱ\\
%       Alveolar tap & \verb|\textfishhookr| & \verb|R| & ɾ \\
%       Retroflex flap & \verb|\textrtailr| & \verb|\:r| & ɽ \\
%       Bilabial voiceless fricative & \verb|\textphi| & \verb|F| & ɸ \\
%       Bilabial voiced fricative & \verb|\textbeta| & \verb|B| & β \\
%       Dental voiceless fricative & \verb|\texttheta| & \verb|T| & θ \\
%       Dental voiced fricative & \verb|\dh|\tnote{\textcolor{red}{§}} &
%       \verb|D| & ð \\
%       Palatal voiceless fricative & \verb|\textesh| & \verb|S| & ʃ \\
%       Palatal voiced fricative & \verb|\textyogh| & \verb|Z| & ʒ \\
%       Retroflex voiceless fricative & \verb|\textrtails| & \verb|\:s| & ʂ \\
%       Retroflex voiceless fricative & \verb|\textrtailz| & \verb|\:z| & ʐ \\
%       Palatal voiced fricative & \verb|\textctj| & \verb|J| & ʝ \\
%       Velar voiced fricative & \verb|\textgamma| & \verb|G| & ɣ \\
%       Uvular voiceless fricative & \verb|\textchi| & \verb|X| & χ \\
%       Uvular voiced fricative & \verb|\textinvscr| & \verb|K| & ʁ \\
%       Pharyngeal voiceless fricative & \verb|\textcrh| & \verb|\*h| & ħ \\
%       Pharyngeal voiced fricative & \verb|\textrevglotstop| & \verb|Q| & ʕ
%       \\
%       Glottal voiced fricative & \verb|\texthth| & \verb|H| & ɦ \\
%       Alveolar voiceless lateral fricative & \verb|\textbeltl| & \verb|\*l|
%       & ɬ \\
%       Alveolar voiced lateral fricative & \verb|\textlyoghlig| & \verb|\*z|
%       & ɮ \\
%       Labiodental approximant & \color{red}{\verb|\labdentapp|} & \verb|V| &
%       ʋ \\
%       Alveolar approximant & \verb|\textturnr| & \verb|\*r| & ɹ \\
%       Retroflex approximant & \verb|\textturnrrtail| & \verb|\:R| & ɻ \\
%       Velar approximant & \verb|\textturnmrleg| & \color{red}{\verb|\*m|} &
%       ɰ \\
%       Retroflex lateral approximant & \verb|\textrtaill| & \verb|\:l| & ɭ \\
%       Palatal lateral approximant & \verb|\textturny| & \verb|L| & ʎ \\
%       Velar lateral approximant & \verb|\textscl| & \verb|\;L| & ʟ \\
%       \midrule
%       \insertTableNotes
%     \end{longtable}
%   \end{ThreePartTable}
%   
%   \subsection{Non-Pulmonic consonants}
%   
%   \begin{longtable}{%
%     l@{\hspace{0.3cm}-\hspace{0.3cm}}%
%     l@{\hspace{0.3cm}-\hspace{0.3cm}}%
%     l@{\hspace{0.3cm}-\hspace{0.3cm}}%
%     >{\ipa[}l<{]}%
%   }%
%     \toprule
%     \endfirsthead
%     \toprule
%     \endhead
%     \bottomrule
%     \endfoot
%     Description & Input A & Input B &\multicolumn{1}{l}{Symbol}\\
%     \midrule
%     Bilabial click & \verb|\textbullseye| & \verb|\!o| & ʘ \\
%     Postalveolar click & \color{red}{\verb|\pstalvclick|} &
%     \color{red}{\verb|!|} & ǃ \\
%     Palatoalveolar click & \verb|\textdoublebarpipe| &
%     \color{red}{\verb|\!c|} & ǂ \\
%     Alveolar lateral click & \verb|\textdoublepipe| & \verb+||+ & ǁ \\
%     Bilabial implosive & \verb|\texthtb| & \verb|\!b| & ɓ \\
%     Dental implosive & \verb|\texthtd| & \verb|\!d| & ɗ \\
%     Palatal implosive & \verb|\texthtbardotlessj| & \verb|\!j| & ʄ \\
%     Velar implosive & \verb|\texthtg| & \verb|\!g| & ɠ \\
%     Uvular implosive & \verb|\texthtscg| & \verb|\!G| & ʛ \\
%   \end{longtable}
%   
%   \subsection{Other symbols}
%   {%
%   \small
%   \begin{longtable}{%
%     >{\footnotesize}l@{\hspace{0.3cm}-\hspace{0.3cm}}%
%     l@{\hspace{0.3cm}-\hspace{0.3cm}}%
%     l@{\hspace{0.3cm}-\hspace{0.3cm}}%
%     >{\ipa[}l<{]}%
%   }%
%     \toprule
%     \multicolumn{1}{l}{Description} & Input A & Input B
%     &\multicolumn{1}{l}{Symbol}\\
%     \midrule
%     Labial-velar voiceless fricative & \verb|\textturnw| & \verb|\*w| & ʍ
%     \\
%     Labial-palatal voiced approximant & \verb|\textturnh| & \verb|4| & ɥ
%     \\
%     Epiglottal voiceless fricative & \verb|\textsch| & \verb|\;H| & ʜ \\
%     Epiglottal voiced fricative/approximant & \verb|\textbarrevglotstop| &
%     \color{red}{\verb|\*1|} & ʢ \\
%     Alveolar voiced lateral flap & \verb|\textturnlonglegr| &
%     \color{red}{\verb|\!r|} & ɺ \\
%     Epiglottal plosive & \verb|\textbarglotstop| & \verb|\*2| & ʡ \\
%     Alveolo-palatal voiceless fricative & \verb|\textctc| & \verb|C| & ɕ \\
%     Alveolo-palatal voiced fricative & \verb|\textctz| &
%     \color{red}{\verb|\*Z|} & ʑ \\
%     Postalveolar-velar voiceless fricative & \verb|\texththeng| &
%     \color{red}{\verb|\*3|} & ɧ \\
%     \bottomrule
%   \end{longtable}
% }
%     
%   \subsection{Vowels}
%   {%
%   \small
%   \begin{longtable}{%
%     l@{\hspace{0.3cm}-\hspace{0.3cm}}%
%     l@{\hspace{0.3cm}-\hspace{0.3cm}}%
%     l@{\hspace{0.3cm}-\hspace{0.3cm}}%
%     >{\ipa[}l<{]}%
%   }%
%     \toprule
%     \endfirsthead
%     \toprule
%     \endhead
%     \bottomrule
%     \endfoot
%     Description & Input A & Input B &\multicolumn{1}{l}{Symbol}\\
%     \midrule
%     Near-close near-front unrounded & \verb|\textsci| & \verb|I| & ɪ \\
%     Near-close near-front rounded & \verb|\textscy| & \verb|Y| & ʏ \\
%     Open-mid front unrounded & \verb|\textepsilon| & \verb|E| & ɛ \\
%     Close central unrounded & \verb|\textbari| & \verb|1| & ɨ \\
%     Close central rounded & \verb|\textbaru| & \verb|0| & ʉ \\
%     Near-close near-back rounded & \verb|\textupsilon| & \verb|U| & ʊ \\
%     Close-mid central unrounded & \verb|\textreve| & \verb|9| & ɘ \\
%     Close-mid central rounded & \verb|\textbaro| & \verb|8| & ɵ \\
%     Mid central vowel & \verb|\textschwa| & \verb|@| & ə \\
%     Open-mid central unrounded & \verb|\textrevepsilon| & \verb|3| & ɜ \\
%     Open-mid central rounded & \verb|\textcloserevepsilon| &
%     \textcolor{red}{\texttt{\&}} & ɞ \\
%     Near-open central & \verb|\textturna| & \verb|5| & ɐ \\
%     Close back unrounded & \verb|\textturnm| & \verb|W| & ɯ \\
%     Close-mid back unrounded & \verb|\textramshorns| & \verb|7| & ɤ \\
%     Open-mid back unrounded & \verb|\textturnv| & \verb|2| & ʌ \\
%     Open-mid back rounded & \verb|\textopeno| & \verb|O| & ɔ \\
%     Open back unrounded & \verb|\textscripta| & \verb|A| & ɑ \\
%     Open back rounded & \verb|\textturnscripta| & \verb|6| & ɒ \\
%   \end{longtable}
% }%
%   
%   \subsection{Independent diacritics}
%   
%   \begin{longtable}{%
%     l@{\hspace{0.3cm}-\hspace{0.3cm}}%
%     l@{\hspace{0.3cm}-\hspace{0.3cm}}%
%     l@{\hspace{0.3cm}-\hspace{0.3cm}}%
%     >{\ipa[}l<{]}%
%   }%
%     \toprule
%     Description & Input A & Input B & \multicolumn{1}{l}{Symbol}\\
%     \midrule
%     Rhoticity & \verb|ə\textrhoticity|\tnotex{1} & \color{red}{\verb|ə\tr|}
%     & ə˞\\
%     No audible release & \verb|n\textcorner| & \color{red}{\verb|n\tc|} & n̚
%     \\
%     Tie bar (standalone) & \color{red}{\verb|\stdlnetiebar{əi}|} &
%     \color{red}{\verb|\ts{əi}|} & ə‿i \\
%     \bottomrule
%   \end{longtable}
%   
%   \subsection{Dependent diacritics}
%   
%   \begin{TableNotes}
%   \item[1] The documentation of package the \pkg{tipa} on p.\ 53
%   describes a shorthand command
%   i.e.~\texttt{\textbackslash\textbar\~{}{\meta{arg}}} which is not
%   functional. I haven't investigated the reasons, but instead provided a new
%   shorthand command for this character.
%   \end{TableNotes}
%   
%   \begin{ThreePartTable}
%     {%
%     \small
%     \begin{longtable}{%
%       l@{\hspace{0.3cm}-\hspace{0.3cm}}%
%       l@{\hspace{0.3cm}-\hspace{0.3cm}}%
%       l@{\hspace{0.3cm}-\hspace{0.3cm}}%
%       >{\ipa[}l<{]}%
%     }%
%       \toprule
%       \endfirsthead
%       \toprule
%       \endhead
%       \bottomrule
%       \endfoot
%       Description & Input A & Input B & \multicolumn{1}{l}{Symbol}\\
%       \midrule
%       Voiceless & \verb|\textsubring{m}| & \verb|\r*m| & m̥\\
%       Voiced & \verb|\textsubwedge{s}| & \verb|\v*s| & s̬\\
%       More rounded & \verb|\textsubrhalfring{w}| & \verb+\|)w+ & w̹\\
%       Less rounded & \verb|\textsublhalfring{w}| & \verb+\|(w+ & w̜\\
%       Advanced & \verb|\textsubplus{ə}| & \verb-\|+e- & ə̟\\
%       Retracted & \verb|\textsubbar{l}| & \verb+\=*l+ & l̠\\
%       Mid-centralized & \verb|\textovercross{e}| & \verb+\|x{e}+ & e̽\\
%       Syllabic & \verb|\textsyllabic{n}| & \verb|\s{n}| & n̩\\
%       Non-syllabic & \verb|\textsubarch{a}| & \color{red}{\verb|\ns{a}|} & a̯
%       \\
%       Breathy voiced & \verb|\textsubumlaut{b}| & \verb|\"*b| & b̤ \\
%       Creaky voiced & \verb|\textsubtilde{ə}| & \verb|\~*@| & ə̰ \\
%       Linguolabial & \verb|\textseagull{l}| & \verb+\|ml+ & l̼ \\
%       Velarized & \verb|\textsuperimposetilde{i}| & \verb+\vl{i}+\tnote{1} & i̴
%       \\
%       Raised & \verb|\textraising{ə}| & \verb+\|'@+ & ə̝ \\
%       Lowered & \verb|\textlowering{ə}| & \verb+\|`@+ & ə̞ \\
%       Advanced tongue root & \verb|\textadvancing{ə}| & \verb+\|<@+ & ə̘ \\
%       Retracted tongue root & \verb|\textretracting{ə}| & \verb+\|>@+ & ə̙ \\
%       Dental & \verb|\textsubbridge{t}| & \verb+\|[t+ & t̪ \\
%       Apical & \verb|\textinvsubbridge{t}| & \verb+\|]t+ & t̺ \\
%       Laminal & \verb|\textsubsquare{s}| & \color{red}{\verb|\lmn{s}|} & s̻ \\
%       Tie bar (top) & \verb|\texttoptiebar{əi}| & \verb|\t{əi}| & ə͡i\\
%       Tie bar (bottom) & \verb|\textbottomtiebar{əi}| & \verb|\t*{əi}| & ə͜i\\
%       \midrule
%       \insertTableNotes
%     \end{longtable}
%   }%
%   \end{ThreePartTable}
% \end{documentation}
% 
% \begin{implementation}
%   \section{Implementation}
%    \begin{macrocode}
%<@@=tipauni>
%<*package>
%    \end{macrocode}
% Provide the package with some description \textit{\&} version specifications.
%    \begin{macrocode}
\ProvidesPackage{tipauni}[2022/02/21 v0.6
Unicode characters with TIPA commands]
%    \end{macrocode}
% Load necessary packages.
%    \begin{macrocode}
\RequirePackage{xparse}
\RequirePackage{fontspec}
\RequirePackage{expkv-def,expkv-opt}
%    \end{macrocode}
% Package options are provided using the package \pkg{expkv}.
%    \begin{macrocode}
\newif\iftipauni@document@font
\def\tipauni@recommended@font{CharisSIL}
\tipauni@document@fonttrue
\let\tipauni@font@options\@empty
\ekvdefinekeys{tipauni}{%
  bool incompatible              = \iftipauni@nontipa@ignore,%
  noval preservefont             = {%
    \let\tipauni@font\tipauni@undefined
    \tipauni@document@fontfalse
  },%
  estore fontspecoptions         = \tipauni@font@options,%
  new noval fontspecoptions      = {%
    \PackageWarningNoLine{tipauni}{%
      As you haven't given any value to package option\MessageBreak
      `fontspecoptions'; it will be ignored. Please\MessageBreak
      provide a valid (list of) fontspec option(s)%
    }%
  },%
  noval resetfontspecoptions     = {%
    \let\tipauni@font@options\@empty
  },%
  noval documentfont             = {%
    \PackageError{tipauni}{No value for the `documentfont' option}{%
      The `documentfont' option has no value. Please provide a valid font-name,%
      \MessageBreak
      otherwise it will be ignored or (if given) the value of another non-empty%
      \MessageBreak
      `documentfont' will be used.%
    }%
    \ifdefined\tipauni@font
      \tipauni@document@fonttrue
    \else
      \tipauni@document@fontfalse
    \fi
  },%
  new code documentfont          = {%
    \edef\tipauni@temp{#1}%
    \ifx\tipauni@temp\@empty
      \PackageError{tipauni}{Empty `documentfont' option}{%
        The `documentfont' option has an empty value. Please provide a valid%
        \MessageBreak
        font-name, otherwise it will be ignored or (if given) the value of%
        \MessageBreak
        another non-empty `documentfont' will be used.%
      }%
      \ifdefined\tipauni@font
        \tipauni@document@fonttrue
      \else
        \tipauni@document@fontfalse
      \fi
    \else
      \tipauni@document@fonttrue
      \let\tipauni@font\tipauni@temp
    \fi
  },%
  noval recommendedfont          = {%
    \tipauni@document@fonttrue
    \let\tipauni@font\tipauni@recommended@font
  }%
}
\ekvoProcessLocalOptions{tipauni}
%    \end{macrocode}
% The following code executes the \verb|\setmainfont| command \& issues a
% warning in a special case.
%    \begin{macrocode}
\iftipauni@document@font
  \ifdefined\tipauni@font
  \else
    \let\tipauni@font\tipauni@recommended@font
  \fi
  \expandafter\setmainfont\expandafter
  [\expandafter{\tipauni@font@options}]{\tipauni@font}
\else
  \ifx\tipauni@font@options\@empty
  \else
    \PackageWarningNoLine{tipauni}{%
      An active font loaded with package fontspec is\MessageBreak
      necessary for `fontspecoptions' to take effect.\MessageBreak
      Please use the `documentfont'/`recommendedfont'\MessageBreak
      option or load a font outside package `tipauni'\MessageBreak
      to make your options active. Currently ignoring:\MessageBreak
      [\tipauni@font@options]%
    }%
  \fi
\fi
%    \end{macrocode}
% The following commands were kindly suggested by Phelype Oleinik for developing
% a conditional which is always global.
%    \begin{macrocode}
\newif\ifnontipa
\protected\def\nontipatrue{\global\let\ifnontipa\iftrue}
\protected\def\nontipafalse{\global\let\ifnontipa\iffalse}
\iftipauni@nontipa@ignore
  \protected\def\nontipatrue{}
\fi
%    \end{macrocode}
% The following commands provide the command \verb|\tipaunicmd| for printing
% characters that are missing in the \pkg{tipa} package.
%    \begin{macrocode}
\NewDocumentCommand\tipaunicmd{ m m }{%
  \NewDocumentCommand#1{ }{\nontipatrue #2}%
}%
%    \end{macrocode}
% As can be seen above, everytime when a command that is absent in the
% \pkg{tipa} package is provided, it sets the conditional \verb|nontipa| to true
% value. At the end of the document this package checks if there is any instance
% of a non-TIPA command. If yes, it generates a warning in the .log file saying
% that this code is now incompatible with the TIPA package and some characters
% might be displayed with a different font. All of this is achieved with the
% help of following commands. The discussion on
% \href{https://topanswers.xyz/tex?q=1735}{this} question helped a lot for fine
% tuning the warning message.
%    \begin{macrocode}
\AtEndDocument{%
  \ifnontipa
  \PackageWarningNoLine{tipauni}{%
    You have used a command(s) that is (are) not\MessageBreak
    supported by package TIPA. Thus your current code is\MessageBreak
    incompatible with that package. These commands are\MessageBreak
    printed with a different Unicode font. Refer to the\MessageBreak
    documentation of package tipauni for a list of\MessageBreak
    commands that are not supported in TIPA%
  }%
  \else
  \fi
}%
%    \end{macrocode}
% From here the code for obtaining Unicode characters with TIPA commands
% starts.
%    \begin{macrocode}
\DeclareTextCommand{\tipaunistar}{TU}[1]{{#1}}
\DeclareTextAccentDefault{\tipaunistar}{TU}
\DeclareTextCommand{\tipaunisemicolon}{TU}[1]{{#1}}
\DeclareTextAccentDefault{\tipaunisemicolon}{TU}
\DeclareTextCommand{\tipaunicolon}{TU}[1]{{#1}}
\DeclareTextAccentDefault{\tipaunicolon}{TU}
\DeclareTextCommand{\tipaunibang}{TU}[1]{{#1}}
\DeclareTextAccentDefault{\tipaunibang}{TU}
\DeclareTextCommand{\tipaunipipe}{TU}[2]{\@pipeshorthands{#1}{#2}}
\DeclareTextAccentDefault{\tipaunipipe}{TU}
%% Pulmonic consonants
%%-- Stops
%%--** Retroflex
\DeclareTextSymbol{\textrtailt}{TU}{"0288}%% ʈ
\DeclareTextComposite{\tipaunicolon}{TU}{t}{"288}
\DeclareTextSymbol{\textrtaild}{TU}{"0256}%% ɖ
\DeclareTextComposite{\tipaunicolon}{TU}{d}{"256}
%%--** Palatal
\DeclareTextSymbol{\textbardotlessj}{TU}{"25F}%% ɟ
\DeclareTextComposite{\tipaunistar}{TU}{j}{"25F}
%%--** Uvular
\DeclareTextSymbol{\textscg}{TU}{"262}%% ɢ
\DeclareTextComposite{\tipaunisemicolon}{TU}{G}{"262}
%%--** Glottal
\DeclareTextSymbol{\textglotstop}{TU}{"294}%% ʔ
%%-- Nasals
%%--** Labiodental
\DeclareTextSymbol{\textltailm}{TU}{"271}%% ɱ
%%--** Retroflex
\DeclareTextSymbol{\textrtailn}{TU}{"273}%% ɳ
\DeclareTextComposite{\tipaunicolon}{TU}{n}{"273}
%%--** Palatal
\DeclareTextSymbol{\textltailn}{TU}{"272}%% ɲ
\DeclareTextComposite{\tipaunistar}{TU}{n}{"272}
%%--** Uvular
\DeclareTextSymbol{\textscn}{TU}{"274}%% ɴ
\DeclareTextComposite{\tipaunisemicolon}{TU}{N}{"274}
%%-- Trills
%%--** Bilabial
\DeclareTextSymbol{\textscb}{TU}{"299}%% ʙ
\DeclareTextComposite{\tipaunisemicolon}{TU}{B}{"299}
%%--** Uvular
\DeclareTextSymbol{\textscr}{TU}{"280}%% ʀ
\DeclareTextComposite{\tipaunisemicolon}{TU}{R}{"280}
%%-- Taps or flaps
%%-- Labiodental flap
\tipaunicmd{\labdentflap}{ⱱ}%% tipauni-exclusive
\DeclareTextComposite{\tipaunistar}{TU}{v}{"2C71}%% tipauni-exclusive
%%--** Alveolar
\DeclareTextSymbol{\textfishhookr}{TU}{"27E}%% ɾ
%%--** Retroflex
\DeclareTextSymbol{\textrtailr}{TU}{"27D}%% ɽ
\DeclareTextComposite{\tipaunicolon}{TU}{r}{"27D}
%%-- Fricatives
%%--** Bilabial
\DeclareTextSymbol{\textphi}{TU}{"278}%% ɸ
\DeclareTextSymbol{\textbeta}{TU}{"3B2}%% β
%%--** Dental
\DeclareTextSymbol{\texttheta}{TU}{"3B8}%% θ
%%--** Postalveolar
\DeclareTextSymbol{\textesh}{TU}{"283}%% ʃ
\DeclareTextSymbol{\textyogh}{TU}{"292}%% ʒ
%%--** Retroflex
\DeclareTextSymbol{\textrtails}{TU}{"282}%% ʂ
\DeclareTextComposite{\tipaunicolon}{TU}{s}{"282}
\DeclareTextSymbol{\textrtailz}{TU}{"290}%% ʐ
\DeclareTextComposite{\tipaunicolon}{TU}{z}{"290}
%%--** Palatal
\DeclareTextSymbol{\textctj}{TU}{"29D}%% ʝ
%%--** Velar
\DeclareTextSymbol{\textgamma}{TU}{"263}%% ɣ
%%--** Uvular
\DeclareTextSymbol{\textchi}{TU}{"3C7}%% χ
\DeclareTextSymbol{\textinvscr}{TU}{"281}%% ʁ
%%--** Pharyngeal
\DeclareTextSymbol{\textcrh}{TU}{"127}%% ħ
\DeclareTextComposite{\tipaunistar}{TU}{h}{"127}
\DeclareTextSymbol{\textrevglotstop}{TU}{"295}%% ʕ
%%--** Glottal
\DeclareTextSymbol{\texthth}{TU}{"266}%% ɦ
%%-- Lateral fricatives
%%--** Alveolar
\DeclareTextSymbol{\textbeltl}{TU}{"26C}%% ɬ
\DeclareTextComposite{\tipaunistar}{TU}{l}{"26C}
\DeclareTextSymbol{\textlyoghlig}{TU}{"26E}%% ɮ
\DeclareTextComposite{\tipaunistar}{TU}{z}{"26E}
%%-- Approximants
%%-- Labiodental approximant
\tipaunicmd{\labdentapp}{ʋ}%% tipauni-exclusive
%%--** Alveolar
\DeclareTextSymbol{\textturnr}{TU}{"279}%% ɹ
\DeclareTextComposite{\tipaunistar}{TU}{r}{"279}
%%--** Retroflex
\DeclareTextSymbol{\textturnrrtail}{TU}{"27B}%% ɻ
\DeclareTextComposite{\tipaunicolon}{TU}{R}{"27B}
%%--** Velar
\DeclareTextSymbol{\textturnmrleg}{TU}{"270}%% ɰ
\DeclareTextComposite{\tipaunistar}{TU}{m}{"270}
%%-- Lateral approximants
%%--** Retroflex
\DeclareTextSymbol{\textrtaill}{TU}{"26D}%% ɭ
\DeclareTextComposite{\tipaunicolon}{TU}{l}{"26D}
\DeclareTextSymbol{\textturny}{TU}{"28E}%% ʎ
\DeclareTextSymbol{\textscl}{TU}{"29F}%% ʟ
\DeclareTextComposite{\tipaunisemicolon}{TU}{L}{"29F}
%% Non-pulmonic consonants
%%-- Clicks
%%--** Bilabial
\DeclareTextSymbol{\textbullseye}{TU}{"298}%% ʘ
\DeclareTextComposite{\tipaunibang}{TU}{o}{"298}
%%-- Postalveolar click
\tipaunicmd{\pstalvclick}{ǃ}%% tipauni-exclusive
%%--** Palato-alveolar
\DeclareTextSymbol{\textdoublebarpipe}{TU}{"1C2}%% ǂ
\DeclareTextComposite{\tipaunibang}{TU}{c}{"1C2}%% tipauni-exclusive
%%--** Alveolar lateral
\DeclareTextSymbol{\textdoublepipe}{TU}{"1C1}% ǁ
%%-- Implosives
%%--** Bilabial
\DeclareTextSymbol{\texthtb}{TU}{"253}%% ɓ
\DeclareTextComposite{\tipaunibang}{TU}{b}{"253}
%%--** Dental
\DeclareTextSymbol{\texthtd}{TU}{"257}%% ɗ
\DeclareTextComposite{\tipaunibang}{TU}{d}{"257}
%%--** Palatal
\DeclareTextSymbol{\texthtbardotlessj}{TU}{"284}%% ʄ
\DeclareTextComposite{\tipaunibang}{TU}{j}{"284}
%%--** Velar
\DeclareTextSymbol{\texthtg}{TU}{"260}%% ɠ
\DeclareTextComposite{\tipaunibang}{TU}{g}{"260}
%%--** Uvular
\DeclareTextSymbol{\texthtscg}{TU}{"29B}%% ʛ
\DeclareTextComposite{\tipaunibang}{TU}{G}{"29B}
%% Other consonants
\DeclareTextSymbol{\textturnw}{TU}{"28D}%% ʍ
\DeclareTextSymbol{\textturnh}{TU}{"265}%% ɥ
\DeclareTextSymbol{\textsch}{TU}{"29C}%% ʜ
\DeclareTextSymbol{\textbarrevglotstop}{TU}{"2A2}%% ʢ
\DeclareTextComposite{\tipaunistar}{TU}{1}{"2A2}%% tipauni-exclusive
\DeclareTextSymbol{\textbarglotstop}{TU}{"2A1}%% ʡ
\DeclareTextComposite{\tipaunistar}{TU}{2}{"2A1}%% tipauni-exclusive
\DeclareTextSymbol{\textctc}{TU}{"255}%% ɕ
\DeclareTextSymbol{\textctz}{TU}{"291}%% ʑ
\DeclareTextComposite{\tipaunistar}{TU}{Z}{"291}%% tipauni-exclusive
\DeclareTextSymbol{\textturnlonglegr}{TU}{"27A}%% ɺ
\DeclareTextComposite{\tipaunibang}{TU}{r}{"27A}%% tipauni-exclusive
\DeclareTextSymbol{\texththeng}{TU}{"267}%% ɧ
\DeclareTextComposite{\tipaunistar}{TU}{3}{"267}%% tipauni-exclusive
\DeclareTextComposite{\tipaunistar}{TU}{k}{"029E}% ʞ
\DeclareTextComposite{\tipaunistar}{TU}{t}{"287}% ʇ
\DeclareTextComposite{\tipaunistar}{TU}{w}{"28D}% ʍ
\DeclareTextComposite{\tipaunisemicolon}{TU}{E}{"1D07}% ᴇ
\DeclareTextComposite{\tipaunisemicolon}{TU}{J}{"1D0A}% ᴊ
\DeclareTextComposite{\tipaunisemicolon}{TU}{A}{"1D00}% ᴀ
\DeclareTextComposite{\tipaunisemicolon}{TU}{H}{"29C}% ʜ
\DeclareTextComposite{\tipaunisemicolon}{TU}{U}{"1D1C}% ᴜ
%% Vowels
\DeclareTextSymbol{\textsci}{TU}{"26A}%% ɪ
\DeclareTextSymbol{\textscy}{TU}{"28F}%% ʏ
\DeclareTextSymbol{\textepsilon}{TU}{"25B}%% ɛ
\DeclareTextSymbol{\textbari}{TU}{"268}%% ɨ
\DeclareTextSymbol{\textbaru}{TU}{"289}%% ʉ
\DeclareTextSymbol{\textupsilon}{TU}{"28A}%% ʊ
\DeclareTextSymbol{\textreve}{TU}{"258}%% ɘ
\DeclareTextSymbol{\textbaro}{TU}{"275}%% ɵ
\DeclareTextSymbol{\textschwa}{TU}{"259}%% ə
\DeclareTextSymbol{\textrevepsilon}{TU}{"25C}%% ɜ
\DeclareTextSymbol{\textcloserevepsilon}{TU}{"25E}%% ɞ
\DeclareTextSymbol{\textturna}{TU}{"250}%% ɐ
\DeclareTextSymbol{\textturnm}{TU}{"26F}%% ɯ
\DeclareTextSymbol{\textramshorns}{TU}{"264}%% ɤ
\DeclareTextSymbol{\textturnv}{TU}{"28C}%% ʌ
\DeclareTextSymbol{\textopeno}{TU}{"254}%% ɔ
\DeclareTextSymbol{\textscripta}{TU}{"251}%% ɑ
\DeclareTextSymbol{\textturnscripta}{TU}{"252}%% ɒ
%% Independent diacritic marks
%%-- Rhoticity
\DeclareTextSymbol{\textrhoticity}{TU}{"2DE}%% ˞
\tipaunicmd{\tr}{˞}
%%-- Tie bar (below)
\DeclareTextCommand{\tipaunitextbottomtiebar}{TU}[2]{#1\symbol{"035C}#2}
\NewDocumentCommand\textbottomtiebar{ m }{\tipaunitextbottomtiebar#1}
%%-- Tie bar (below) - standalone character
\DeclareTextCommand{\tipaunits}{TU}[2]{#1\symbol{"203F}#2}
\NewDocumentCommand\stdlnetiebar{ m }{\tipaunits#1}
\NewDocumentCommand\ts{ m }{\tipaunits#1}
%    \end{macrocode}
% The following commands redefine the TIPA shorthand input for dependent
% diacritics.
%    \begin{macrocode}
%% No audible release
\DeclareTextSymbol{\textcorner}{TU}{"31A}%% ̚
\tipaunicmd{\tc}{̚}
%% Syllabic
\DeclareUnicodeAccent{\textsyllabic}{TU}{"0329}%% ̩
\NewDocumentCommand\s{ m }{\textsyllabic{\textipa{#1}}\nontipatrue}
%% Non-syllabic
\DeclareUnicodeAccent{\textsubarch}{TU}{"032F}%% ̯
\NewDocumentCommand\ns{ m }{\textsubarch{\textipa{#1}}}
%% Laminal
\DeclareUnicodeAccent{\textsubsquare}{TU}{"033B}%% ̻
\NewDocumentCommand\lmn{ m }{\textsubsquare{\textipa{#1}}\nontipatrue}
%% Other complex shorthand commands
\DeclareTextCommand{\@pipeshorthands}{TU}[2]{%
  \if #1[\tipaunisubbridge{\textipa{#2}}\else
  \if #1]\tipauniinvsubbridge{\textipa{#2}}\else
  \if #1(\tipaunisublhalfring{\textipa{#2}}\else
  \if #1)\tipaunisubrhalfring{\textipa{#2}}\else
  \if #1+\tipaunisubplus{\textipa{#2}}\else
  \if #1'\tipauniraising{\textipa{#2}}\else
  \if #1`\tipaunilowering{\textipa{#2}}\else
  \if #1<\tipauniadvancing{\textipa{#2}}\else
  \if #1>\tipauniretracting{\textipa{#2}}\else
  \if #1x\tipauniovercross{\textipa{#2}}\else
  \if #1m\tipauniseagull{\textipa{#2}}\else #1#2%
  \fi\fi\fi\fi\fi\fi\fi\fi\fi\fi\fi
}%
\DeclareTextAccentDefault{\@pipeshorthands}{TU}
\ExplSyntaxOn
\NewDocumentCommand \super { m }
  {%
    \str_case:nn {#1}
      {
        {h} {ʰ}
        {w} {ʷ}
        {j} {ʲ}
        {ɣ} {ˠ}
        {ʕ} {ˤ}
        {n} {ⁿ}
        {l} {ˡ}
      }
  }
%    \end{macrocode}
% We need a few variables.
%    \begin{macrocode}
\seq_new:N \l_@@_remove_from_accents_seq
\tl_new:N \l_@@_textipa_tl
%    \end{macrocode}
% Since we have a non-standard syntax for some of the accent macros (having a
% starred variant) we have to exclude them from the list of accents which get
% special handling from \cs{text_expand:n} or else the stars will be surrounded
% by braces and not found by \cs{@ifstar}. This sequence stores the problematic
% accents which need to be excluded.
%    \begin{macrocode}
\seq_set_split:Nnn \l_@@_remove_from_accents_seq { , }
  { \~, \=, \", \r, \v, \t }
%    \end{macrocode}
% We only search for the macros not contained in \pkg{tipa} if this is requested
% by the user.
%    \begin{macrocode}
\iftipauni@nontipa@ignore
  \cs_new_eq:NN \@@_nontipa_search: \prg_do_nothing:
  \cs_new_eq:NN \@@_nontipa_search:n \use_none:n
\else
  \cs_new_protected:Npn \@@_nontipa_search:
    {
      \unless\ifnontipa
        \exp_args:No \@@_nontipa_search:n \l_@@_textipa_tl
      \fi
    }
%    \end{macrocode}
% The warning thrown at the end of the run doesn't list which symbols were
% incompatible, so we stop searching when we found one incompatible symbol. We
% use \cs{str_if_in:nnT} so that groups don't affect the search. The alternative
% would be the much slower \cs{regex_match:nnT}.
%    \begin{macrocode}
  \cs_new_protected:Npn \@@_nontipa_search:n #1
    {
      \str_if_in:nnT {#1} {ⱱ} { \nontipatrue \use_none_delimit_by_s_stop:w }
      \str_if_in:nnT {#1} {ɞ} { \nontipatrue \use_none_delimit_by_s_stop:w }
      \str_if_in:nnT {#1} {ɰ} { \nontipatrue \use_none_delimit_by_s_stop:w }
      \str_if_in:nnT {#1} {ǃ} { \nontipatrue \use_none_delimit_by_s_stop:w }
      \str_if_in:nnT {#1} {ǂ} { \nontipatrue \use_none_delimit_by_s_stop:w }
      \str_if_in:nnT {#1} {ʢ} { \nontipatrue \use_none_delimit_by_s_stop:w }
      \str_if_in:nnT {#1} {ɺ} { \nontipatrue \use_none_delimit_by_s_stop:w }
      \str_if_in:nnT {#1} {ʑ} { \nontipatrue \use_none_delimit_by_s_stop:w }
      \str_if_in:nnT {#1} {ɧ} { \nontipatrue \use_none_delimit_by_s_stop:w }
      \use_none:nnn \s_stop
      \use:n { \cs_gset_eq:NN \@@_nontipa_search: \prg_do_nothing: }
    }
\fi
%    \end{macrocode}
% The user facing macro needs to make some local settings.
%    \begin{macrocode}
\NewDocumentCommand\textipa{ +m }%
  {
    \begingroup
      \let\*\tipaunistar
      \let\:\tipaunicolon
      \let\;\tipaunisemicolon
      \let\!\tipaunibang
      \let\|\tipaunipipe
      \let\t\tipaunit
%    \end{macrocode}
% Here we locally remove the accents which have a starred variant from the
% exclusion list of \cs{text_expand:n}.
%    \begin{macrocode}
      \seq_map_inline:Nn \l_@@_remove_from_accents_seq
      { \tl_remove_all:Nn \l_text_accents_tl {##1} }
%    \end{macrocode}
% Here is the real replacement. First we carefully expand the input using
% \cs{text_expand:n}, then we'll loop over the token list with a construct
% similar to the \pkg{expl3} internal \cs{__tl_act:NNNn} (hence the name
% \cs{@@_act:n}). Since the replacement code depends on the correct string
% representation of things we explicitly set the \cs{escapechar} to a backslash
% and reset it afterwards.
%    \begin{macrocode}
      \tl_set:Nx \l_@@_textipa_tl { \text_expand:n {#1} }
      \use:x
        {
          \exp_not:n
            {
              \escapechar = `\\
              \tl_set:Nx \l_@@_textipa_tl
                { \exp_args:No \@@_act:n \l_@@_textipa_tl }
              \escapechar =
            }
          \int_use:N \escapechar
        }
      \exp_stop_f:
      \l_@@_textipa_tl
%    \end{macrocode}
% Finally test for \pkg{tipa}-incompatible symbols.
%    \begin{macrocode}
      \@@_nontipa_search:
    \endgroup
  }
%    \end{macrocode}
% We'll need a scan mark for the |act| code, so we initialise it.
%    \begin{macrocode}
\scan_new:N \s_@@_act_stop
\cs_new:Npn \@@_use_none_delimit_by_act_stop:w #1 \s_@@_act_stop {}
%    \end{macrocode}
% The act macro uses a bit of low level coding to get good performance. It
% implements a loop which does the replacement of single letters, bigraphs or
% combinations of one macro and one letter (which might be separated from the
% macro by multiple spaces, or contained in braces as an argument, or both).
% \cs{_@@_act:n} should be used inside an |x|-expansion context. The only
% token which is not allowed to be ever part of the argument for
% \cs{_@@_act:n} is \cs{s_@@_act_stop}.
%    \begin{macrocode}
\cs_new:Npn \@@_act:n #1 { \@@_act_loop:w #1 {\s_@@_act_stop} . \s_@@_act_stop }
%    \end{macrocode}
% Just to have an easier to spot output we use a private copy of \cs{exp_not:n},
% meaning we leave the output in the input stream (but since the whole loop
% should be used in an |x|-expansion context this is fine, anything expecting
% arguments will get those later).
%    \begin{macrocode}
\cs_new_eq:NN \@@_act_output:n \exp_not:n
%    \end{macrocode}
% The current loop status is indicated by one macro -- \cs{use_none:nn} for the
% initial state; \cs{@@_act_macro:nn} if a macro being the start of a
% bigraph was encountered; or \cs{@@_act_char:nn} if a character being
% the start of a bigraph or a monograph was encountered -- followed by one
% |n|-type argument containing information depending on the state. So for
% \cs{@@_act_loop:Nnw} the arguments are: |#1| is the status macro, |#2|
% more status information, |#3| the remainder of the argument list. It'll branch
% off depending on the first token in |#3| being either a space, a group
% (anything started by a character of category 1, typically an opening brace)
% and anything else (a valid |N|-type argument).
%    \begin{macrocode}
\cs_new:Npn \@@_act_loop:w { \@@_act_loop:Nnw \use_none:nn {} }
\cs_new:Npn \@@_act_loop:Nnw #1#2#3 \s_@@_act_stop
  {
    \@@_act_if_head_is_space:nTF {#3}
      { \@@_act_space:Nnw }
      {
        \tl_if_head_is_group:nTF {#3}
          \@@_act_group:Nnn
          \@@_act_normal:NnN
      }
      #1 {#2} #3 \s_@@_act_stop
  }
%    \end{macrocode}
% We'll need a few explicit space tokens in the following definitions, so we use
% a temporary macro to set this up.
%    \begin{macrocode}
\group_begin:
  \cs_set:Npn \@@_tmp:n #1
    {
%    \end{macrocode}
% Quick test for a leading space using the |stop|-marker.
%    \begin{macrocode}
      \cs_new:Npn \@@_act_if_head_is_space:nTF ##1
        {
          \@@_act_if_head_is_space:w
            \s_@@_act_stop ##1 \s_@@_act_stop
            \@@_act_if_head_is_space_true:w
            \s_@@_act_stop #1 \s_@@_act_stop
            \use_ii:nn
        }
      \cs_new:Npn \@@_act_if_head_is_space:w
          ##1 \s_@@_act_stop #1 ##2 \s_@@_act_stop
        {}
      \cs_new:Npn \@@_act_if_head_is_space_true:w
          \s_@@_act_stop #1 \s_@@_act_stop \use_ii:nn ##1 ##2
        {##1}
%    \end{macrocode}
% The case with a space needs a bit of special treatment if the current state is
% |macro|. In that case we defer the decision and save the space inside the
% additional status |n|-type argument. That consists of two nested |n|-type
% arguments, the first being the control sequence under test and the second
% tokens which need to be reinserted should the replacement fail.
%    \begin{macrocode}
      \cs_new:Npn \@@_act_space:Nnw ##1 ##2 #1
        {
          \token_if_eq_meaning:NNTF ##1 \@@_act_macro:nn
            { \@@_act_space_macro:nn ##2 }
            {
              \cs_if_exist_use:cTF { @@_act_char_mono_ \tl_to_str:n { ##2 } : }
                { \@@_act_output:n {~} }
                { \@@_act_output:n { ##2 ~ } }
              \@@_act_loop:w
            }
        }
    }
  \@@_tmp:n {~}
\group_end:
%    \end{macrocode}
% This macro just works as a helper to get the two separate arguments in the
% status and append a space token to the deferred list.
%    \begin{macrocode}
\cs_new:Npn \@@_act_space_macro:nn #1#2
  { \@@_act_loop:Nnw \@@_act_macro:nn { {#1} { #2 ~ } } }
%    \end{macrocode}
% The end marker is a single \cs{s_@@_act_stop} in a group, hence we need to
% test for that. If this isn't the end we need to handle two different cases, if
% the current state is a macro the group might actually be the correct argument
% to that macro, else a bigraph shouldn't contain a group, so we output what
% ever the current status deferred information was (either as a |mono|
% replacement or just as is) and act on the group.
%    \begin{macrocode}
\cs_new:Npn \@@_act_group:Nnn #1#2#3
  {
    \@@_use_none_delimit_by_act_stop:w
      #3 \@@_act_final:Nnw #1 {#2} \s_@@_act_stop
    \token_if_eq_meaning:NNTF #1 \@@_act_macro:nn
      { \@@_act_group_macro:nn {#2} {#3} }
      {
        \cs_if_exist_use:cF { @@_act_char_mono_ \tl_to_str:n { #2 } : }
          { \@@_act_output:n {#2} }
        { \@@_act_loop:w #3 {\s_@@_act_stop} . \s_@@_act_stop }
      }
    \@@_act_loop:w
  }
%    \end{macrocode}
% The final code needs to remove the remainder of the loop and reinsert whatever
% was deferred, unless that happens to be one of the single tokens to be
% replaced in which case we do the replacement.
%    \begin{macrocode}
\cs_new:Npn \@@_act_final:Nnw
    #1#2 \s_@@_act_stop #3 \@@_act_loop:w #4 \s_@@_act_stop
  {
    \cs_if_exist_use:cF { @@_act_char_mono_ \tl_to_str:n { #2 } : }
      {
        \token_if_eq_meaning:NNTF #1 \@@_act_macro:nn
          { \exp_args:No \@@_act_output:n { \use:nn #2 } }
          { \@@_act_output:n {#2} }
      }
  }
%    \end{macrocode}
% The normal case is pretty straight forward, if the current state is the
% initial state, check whether the token under test is the start of one of the
% replacement sequences (this would be stored inside a macro), if it isn't
% directly output it. If it is the code in that macro will do the right thing.
% If the current state is something else the state macro will handle things, so
% we just call it with the arguments it should get.
%    \begin{macrocode}
\cs_new:Npn \@@_act_normal:NnN #1#2#3
  {
    \token_if_eq_meaning:NNTF #1 \use_none:nn
      {
        \cs_if_exist_use:cF { @@_act_normal_ \tl_to_str:n {#3} :nw }
          { \@@_act_output:n }
          {#3}
      }
      { #1 {#2} {#3} }
    \@@_act_loop:w
  }
%    \end{macrocode}
% The code that should handle a macro will test whether a replacement occurs
% (stored inside a macro), if it does calls that macro, else outputs the
% deferred information. Since the |macro| state will have two groups in the
% deferred state we need to handle that using \cs{use_i:nn} to get the macro
% name, and \cs{use:nn} to just unbrace the macro and the other deferred tokens
% (which might be any number of spaces).
%    \begin{macrocode}
\cs_new:Npn \@@_act_macro:nn #1#2
  {
    \cs_if_exist_use:cF
      { @@_act_macro_ \exp_args:No \tl_to_str:n { \use_i:nn #1 / #2 } : }
      { \exp_args:No \@@_act_output:n { \use:nn #1 #2 } }
  }
%    \end{macrocode}
% If the |macro| state was active when a group was encountered we need to check
% whether that group was the correct argument for the replacement, if it was do
% the replacement, else reinsert the macro and deferred tokens and loop through
% the group.
%    \begin{macrocode}
\cs_new:Npn \@@_act_group_macro:nn #1#2
  {
    \cs_if_exist_use:cF
      { @@_act_macro_ \exp_args:No \tl_to_str:n { \use_i:nn #1 / #2 } : }
      {
        \exp_args:No \@@_act_output:n { \use:nn #1 }
        { \@@_act_loop:w #2 {\s_@@_act_stop} . \s_@@_act_stop }
      }
  }
%    \end{macrocode}
% If the |char| state is active we need to check whether a bigraph was found, if
% so do that replacement, else test whether the single char is to be replaced
% without being part of a bigraph, if not output it unchanged. If there was no
% bigraph the newly found character might be the start of a mono- or bigraph so
% we pass it back to the |N|-type act to evaluate that, in which case we need
% to remove a trailing \cs{@@_act_loop:w} which would start the next loop which
% we do using \cs{use_i:nn}.
%    \begin{macrocode}
\cs_new:Npn \@@_act_char:nn #1#2
  {
    \cs_if_exist_use:cF { @@_act_char_bigraph_ \tl_to_str:n { #1/#2 } : }
      {
        \cs_if_exist_use:cF { @@_act_char_mono_ \tl_to_str:n { #1 } : }
          { \@@_act_output:n {#1} }
        \use_i:nn { \@@_act_normal:NnN \use_none:nn {} {#2} }
      }
  }
%    \end{macrocode}
% To register new bigraphs we use the following functions. The information is
% always stored inside two macros, the first is called
% |@@_act_normal_|\meta{first token}|:nw| and will set the correct state, the
% second is named dependent on the state it should be found in, either
% |@@_act_macro_|\meta{first token}|/|\meta{second token}|:|,
% |@@_act_char_bigraph_|\meta{first token}|/|\meta{second token}|:|, or
% |@@_act_char_mono_|\meta{token}|:|. Those second macros contain the actual
% replacement.
%
% The first macro is a wrapper which just inputs \cs{@@_new_macro_bigraph:NNn}
% for the two macros in |#1| and |#2|.
%    \begin{macrocode}
\cs_new_protected:Npn \@@_new_macro_pair_bigraph:NNNn #1#2#3#4
  {
    \@@_new_macro_bigraph:NNn #1 #3 {#4}
    \@@_new_macro_bigraph:NNn #2 #3 {#4}
  }
%    \end{macrocode}
% A macro bigraph should consist of a control sequence followed by another
% token. It'll switch the state to |macro|.
%    \begin{macrocode}
\cs_new_protected:Npn \@@_new_macro_bigraph:NNn #1#2#3
  {
    \cs_gset:cpn { @@_act_normal_ \tl_to_str:n {#1} :nw } ##1##2 \@@_act_loop:w
      { \@@_act_loop:Nnw \@@_act_macro:nn { {##1} {} } }
    \cs_gset:cpn { @@_act_macro_ \tl_to_str:n { #1/#2 } : }
      { \@@_act_output:n {#3} }
  }
%    \end{macrocode}
% A character bigraph should consist of two characters. It'll switch the state
% to |char|.
%    \begin{macrocode}
\cs_new_protected:Npn \@@_new_char_bigraph:NNn #1#2#3
  {
    \cs_gset:cpn { @@_act_normal_ \tl_to_str:n {#1} :nw } ##1##2 \@@_act_loop:w
      { \@@_act_loop:Nnw \@@_act_char:nn {##1} }
    \cs_gset:cpn { @@_act_char_bigraph_ \tl_to_str:n { #1/#2 } : }
      { \@@_act_output:n {#3} }
  }
%    \end{macrocode}
% And lastly the single replacements which should just be a single character.
% These will also switch the state to |char| (as they might be also a start of
% a bigraph).
%    \begin{macrocode}
\cs_new_protected:Npn \@@_new_char_mono:Nn #1#2
  {
    \cs_gset:cpn { @@_act_normal_ \tl_to_str:n {#1} :nw } ##1##2 \@@_act_loop:w
      { \@@_act_loop:Nnw \@@_act_char:nn {##1} }
    \cs_gset:cpn { @@_act_char_mono_ \tl_to_str:n {#1} : }
      { \@@_act_output:n {#2} }
  }
%    \end{macrocode}
% Everything that's now missing is defining the macros storing the replacements
% as described above.
%    \begin{macrocode}
\@@_new_macro_pair_bigraph:NNNn \! \tipaunibang      G {ʛ}
\@@_new_macro_pair_bigraph:NNNn \! \tipaunibang      g {ɠ}
\@@_new_macro_pair_bigraph:NNNn \* \tipaunistar      1 {ʢ}
\@@_new_macro_pair_bigraph:NNNn \* \tipaunistar      2 {ʡ}
\@@_new_macro_pair_bigraph:NNNn \* \tipaunistar      3 {ɧ}
\@@_new_macro_pair_bigraph:NNNn \* \tipaunistar      Z {ʑ}
\@@_new_macro_pair_bigraph:NNNn \: \tipaunicolon     R {ɻ}
\@@_new_macro_pair_bigraph:NNNn \; \tipaunisemicolon A {ᴀ}
\@@_new_macro_pair_bigraph:NNNn \; \tipaunisemicolon G {ɢ}
\@@_new_macro_pair_bigraph:NNNn \; \tipaunisemicolon B {ʙ}
\@@_new_macro_pair_bigraph:NNNn \; \tipaunisemicolon E {ᴇ}
\@@_new_macro_pair_bigraph:NNNn \; \tipaunisemicolon H {ʜ}
\@@_new_macro_pair_bigraph:NNNn \; \tipaunisemicolon J {ᴊ}
\@@_new_macro_pair_bigraph:NNNn \; \tipaunisemicolon L {ʟ}
\@@_new_macro_pair_bigraph:NNNn \; \tipaunisemicolon N {ɴ}
\@@_new_macro_pair_bigraph:NNNn \; \tipaunisemicolon R {ʀ}
\@@_new_macro_pair_bigraph:NNNn \; \tipaunisemicolon U {ᴜ}
\@@_new_char_bigraph:NNn || {ǁ}
\@@_new_char_bigraph:NNn "" {ˌ}
\clist_map_inline:nn
  {
    {5}{ɐ}, {A}{ɑ}, {6}{ɒ}, {2}{ʌ}, {B}{β}, {C}{ɕ}, {D}{ð}, {9}{ɘ}, {E}{ɛ},
    {3}{ɜ}, {g}{ɡ}, {G}{ɣ}, {7}{ɤ}, {H}{ɦ}, {L}{ʎ}, {4}{ɥ}, {1}{ɨ}, {I}{ɪ},
    {J}{ʝ}, {M}{ɱ}, {W}{ɯ}, {N}{ŋ}, {8}{ɵ}, {O}{ɔ}, {F}{ɸ}, {R}{ɾ}, {K}{ʁ},
    {S}{ʃ}, {T}{θ}, {0}{ʉ}, {U}{ʊ}, {V}{ʋ}, {X}{χ}, {Y}{ʏ}, {Z}{ʒ},
    {P}{ʔ}, {Q}{ʕ},
    {:}{ː}, {;}{ˑ}, {!}{ǃ}, {&}{ɞ}, {@}{ə}, {|}{ǀ}, {"}{ˈ}
  }
  { \@@_new_char_mono:Nn #1 }
\ExplSyntaxOff
%    \end{macrocode}
% The following commands define the long commands for dependent diacritics.
%    \begin{macrocode}
%%-- Dependent diacritic marks
%%--** Voiceless
\DeclareUnicodeAccent{\textsubring}{TU}{"0325}%% ̥
\DeclareTextCommand{\r}{TU}{\tipaunisubring}
\protected\def\tipaunisubring{\@ifstar{\TipaUniSubRing}{\TipaUniSubRing}}
\NewDocumentCommand\TipaUniSubRing{ m }{\textsubring{\textipa{#1}}}
%%--** Voiced
\DeclareUnicodeAccent{\textsubwedge}{TU}{"032C}%% ̬
\DeclareTextCommand{\v}{TU}{\tipaunisubwedge}
\protected\def\tipaunisubwedge{\@ifstar{\TipaUniSubWedge}{\TipaUniSubWedge}}
\NewDocumentCommand\TipaUniSubWedge{ m }{\textsubwedge{\textipa{#1}}}
%%--** More rounded
\DeclareUnicodeAccent{\textsubrhalfring}{TU}{"0339}%% ̹
\NewDocumentCommand\tipaunisubrhalfring{ m }{\textsubrhalfring{\textipa{#1}}}
%%--** Less rounded
\DeclareUnicodeAccent{\textsublhalfring}{TU}{"031C}%% ̜
\NewDocumentCommand\tipaunisublhalfring{ m }{\textsublhalfring{\textipa{#1}}}
%%--** Advanced
\DeclareUnicodeAccent{\textsubplus}{TU}{"031F}%% ̟
\NewDocumentCommand\tipaunisubplus{ m }{\textsubplus{\textipa{#1}}}
%%--** Retracted
\DeclareUnicodeAccent{\textsubbar}{TU}{"0320}%% ̠
\DeclareTextCommand{\=}{TU}{\tipaunisubbar}
\protected\def\tipaunisubbar{\@ifstar{\TipaUniSubBar}{\TipaUniSubBar}}
\NewDocumentCommand\TipaUniSubBar{ m }{\textsubbar{\textipa{#1}}}
%%--** Mid-centralized
\DeclareUnicodeAccent{\textovercross}{TU}{"033D}%% ̽
\NewDocumentCommand\tipauniovercross{ m }{\textovercross{\textipa{#1}}}
%%--** Breathy voiced
\DeclareUnicodeAccent{\textsubumlaut}{TU}{"0324}%% ̤
\DeclareTextCommand{\"}{TU}{\tipaunisubumlaut}
\protected\def\tipaunisubumlaut{\@ifstar{\TipaUniSubUmlaut}{\TipaUniSubUmlaut}}
\NewDocumentCommand\TipaUniSubUmlaut{ m }{\textsubumlaut{\textipa{#1}}}
%%--** Creaky voiced
\DeclareUnicodeAccent{\textsubtilde}{TU}{"0330}%% ̰
\DeclareTextCommand{\~}{TU}{\tipaunisubtilde}
\protected\def\tipaunisubtilde{\@ifstar{\TipaUniSubTilde}{\TipaUniSubTilde}}
\NewDocumentCommand\TipaUniSubTilde{ m }{\textsubtilde{\textipa{#1}}}
%%--** Linguo-labialized
\DeclareUnicodeAccent{\textseagull}{TU}{"033C}%% ̼
\NewDocumentCommand\tipauniseagull{ m }{\textseagull{#1}}
%%--** Velarized
\DeclareUnicodeAccent{\textsuperimposetilde}{TU}{"0334}%% ̴
\NewDocumentCommand\tipaunisuperimposetilde{ m }{\textsuperimposetilde{#1}}
\NewDocumentCommand\vl{ m }{\textsuperimposetilde{\textipa{#1}}}
%%--** Raised
\DeclareUnicodeAccent{\textraising}{TU}{"031D}%% ̝
\NewDocumentCommand{\tipauniraising}{m}{\textraising{#1}}
%%--** Lowered
\DeclareUnicodeAccent{\textlowering}{TU}{"031E}%% ̞
\NewDocumentCommand{\tipaunilowering}{m}{\textlowering{#1}}
%%%--** Advanced tongue root
\DeclareUnicodeAccent{\textadvancing}{TU}{"0318}%% ̘
\NewDocumentCommand{\tipauniadvancing}{m}{\textadvancing{#1}}
%%--** Retracted tongue root
\DeclareUnicodeAccent{\textretracting}{TU}{"0319}%% ̙
\NewDocumentCommand{\tipauniretracting}{m}{\textretracting{#1}}
%%--** Dental
\DeclareUnicodeAccent{\textsubbridge}{TU}{"032A}%% ̪
\NewDocumentCommand\tipaunisubbridge{ m }{\textsubbridge{#1}}
%%--** Apical
\DeclareUnicodeAccent{\textinvsubbridge}{TU}{"033A}%% ̺
\NewDocumentCommand\tipauniinvsubbridge{ m }{\textinvsubbridge{#1}}
%%--** Tie bar (above)
\DeclareTextCommand{\tipaunitexttoptiebar}{TU}[2]{#1\symbol{"0361}#2}
\NewDocumentCommand\texttoptiebar{ m }{\tipaunitexttoptiebar#1}
\DeclareTextCommand{\tipaunit}{TU}{\@ifstar{\textbottomtiebar}{\texttoptiebar}}
%    \end{macrocode}
% Lastly the code for environment \texttt{IPA} is provided by the following.
%    \begin{macrocode}
\NewDocumentEnvironment{IPA}{+b}{\textipa{#1}}{}
%    \end{macrocode}
%    \begin{macrocode}
%</package>      
%    \end{macrocode}
% \end{implementation}
%
% \pagebreak
%
% \input{gfdl-tex.tex}
% 
% \Finale
